% References Figure~\ref{fig:highlights:lsstcam}, defined in the summary figure near the top of researchStatement.tex.
\textbf{Astronomical Instrumentation Development:}
My recent work has spanned both photometric and spectroscopic instruments.
This began in 2020, with the \LLAMAS, where I validated the optical performance of the diffraction gratings and collimators for the \LLAMAS\ spectrographs, assembled and verified optical tolerance for the 24 spectrograph cameras, and integrated all 2,400 AR-coated fibers into the focal plane with zero fiber loss, before it was commissioned on the Magellan Baade telescope in 2024. % LLAMAS
% After integrating and testing \LLAMAS, I shifted my focus from integral-field spectrographs to multi-object spectrographs. Starting in 2023, while at the University of Michigan-Ann Arbor, I tested a common failure mode on \DESI\ robotic positioners, where I found that reorienting the robotic positioner to an inverted physical configuration could eliminate the failure mode. This indicated that gravity-sag has a significant influence on the presence and recovery of these failures \citep{Wuthrich2026Positioners}. % DESI testing
This expertise with spectroscopic instruments translated effectively to the transcendent photometric imager of the next decade, the \LSSTCam. Shortly after beginning my PhD at the University of Zürich, I moved to Chile in 2024 to serve as an \LSSTCam\ Summit Support Scientist, where \textbf{I led \LSSTCam\ focal-plane commissioning} on Cerro Pach\'on as a principal component of the in-kind contribution from Switzerland to the Rubin Project. In this role, I tested the CCDs in both the summit clean-room and on the Simonyi Survey Telescope, and optimized the sensor performance. This played a leading role in achieving \LSSTCam\ first photon in April 2024.
% This work culminated in an invited proceeding to the SPIE Astronomical Telescopes + Instrumentation conference in July 2026, an impactful astronomical instrumentation conference held biennially \citep{MacBride2026LSSTCam}. % LSSTCam testing
% Additionally, I characterized CCD defects in the \LSSTCam\ focal plane, and developed a custom image-processing tool to monitor stability of defects during \LSSTCam operation. % LSSTCam

% References Figure~\ref{fig:highlights:s251112cm_search}, defined in the summary figure near the top of researchStatement.tex.
\textbf{Target-of-Opportunity Searches:}
Since 2023, I have led twenty observing campaigns using different telescopes and facilities, with the goal of detecting \EM\ counterparts to \GW\ merger events. 
A variety of multi-messenger analyses are enabled by the detection of an \EM\ counterpart to a \GW\ merger event, and my particular focus is to combine the redshift information from the host galaxy of an \EM\ counterpart with the luminosity distance inferred from the \GW\ source. With this combination, it is possible to constrain the local rate of accelerated expansion of the universe, known as the Hubble Constant ($H_0$) in the standard model of cosmology, $\Lambda$CDM. This technique is known as the \textit{bright siren} method.

Since 2023, I have have been co-leading the \DESGW\ program, conducting \DECam\ observations, processing raw image data from sixteen \ToO\ campaigns to calibrated products, and analyzing processed data in response to \LVK\ Observing Run 4 (O4) triggers.
During a recent observing campaign in 2025, I led a collaboration between \DESGW\ and the \JGEM\ when a well-localized gravitational wave merger candidate was accessible to both Northern and Southern hemisphere instruments. In this collaboration, I obtained multi-epoch photometry using \DECam, processed the obtained data, and cross-matched this data with spectroscopic observations using \PFS, a recently commissioned multi-object spectrograph at the Subaru Telescope. This analysis enabled the first scientific result using the \PFS.
% This data serves a dual purpose, providing valuable transient characterization information, and also serving as the foundation for a measurement of \Ho\ using the dark siren method with spectroscopic data (see Figure~\ref{fig:pfs_coverage}).
% In addition to maintaining the day-to-day operations of the DESGW program during the IGWN O4 observing run, I improved the gravitational-wave alert-processing pipeline to provide high quality, low-latency information, to inform \ToO\ triggering decisions. Alongside improvements to the alert-processing pipeline, I augmented DESGW observing strategy by implementing changes to re-weight specific sky areas based on a variety of environmental and astrophysical conditions \TODO{Cite the MI paper, even if in-prep, plus a concrete metric from the paper}.

\textbf{Since 2025, I have been leading the \Rubin\ \ToO\ observing program as the \Rubin\ \ToO\ observer}. In this role, I have led observing responses to one interstellar object, one high-energy neutrino event (later reclassified as a terrestrial signal), and two gravitational-wave merger candidates, including the deepest, wide-area \GW\ \ToO\ search to date ($m_{g,5\sigma}\sim24.8$ across $\sim850\deg^2$). % Rubin ToO
% The process to build, test, and deploy an automated \ToO\ system for \Rubin\ is reflected in the recent proceedings of SPIE Astronomical Telescopes + Instrumentation 2026 \citep{MacBride2026RubinToO}, closely following the community-defined science case for the program \citep{Andreoni2024RubinToO}.

% In addition to building the \Rubin\ \ToO\ system, I have led extensive testing of the \Rubin\ \ToO\ infrastructure with the \IGWN\ and \IceCube\ collaborations, including two mock-data-challenges to assess the average latency of a \Rubin\ \ToO\ response, the accuracy of our alert processing pipeline, and the sensitivity of simulated \Rubin\ observations to optical counterparts \citep{MacBride2025MDC}\TODO{Add a metric from the mock-data-challenge}.


% References Figure~\ref{fig:highlights:delve_result}, defined in the summary figure near the top of researchStatement.tex.
\textbf{Standard Siren Constraints on the Hubble Constant:}
While recent observing campaigns have not yet yielded an \EM\ counterpart to a \GW\ merger, I have retained a focus on constraining cosmology using \EM\ and \GW\ data, employing the \textit{dark siren} method to measure the local expansion rate of the universe, \Ho. 
The dark siren method infers redshift information of \GW\ events using a catalog of potential host galaxies, which is lower precision than the bright siren method, but can be applied to all GW events with galaxy catalog coverage.
My work began by supporting the preparation of \DES\ data, resulting in the first dark siren constraint using a deep photometric catalog to jointly infer \GW\ population parameters and cosmological parameters. \DES\ was accepted as the fiducial result for the most recent \LVK\ analysis, and demonstrated a 22\% improvement in constraining \Ho\ compared to previous analyses. I have since expanded on the \DES\ preparation, leading a dark-siren analysis using \DELVE, where initial results show that when photometric redshift biases are properly treated, \textbf{\DELVE\ outperforms \DES\ using dark sirens}, constraining \Ho\ to 11\% when combined with the bright siren GW170817. % DES+DELVE

